% Luc Destombe --- licence CC BY-NC-SA 4.0
% https://creativecommons.org/licenses/by-nc-sa/4.0/deed.fr
% Fichier autonome : compiler avec pdflatex, deux fois (signets, nombre de pages).
\documentclass[11pt,a4paper]{article}
\makeatletter
\newif\ifeleve
\elevefalse
\elevetrue

\newif\iflycee@palatino
\lycee@palatinotrue

\RequirePackage[utf8]{inputenc}
\RequirePackage[T1]{fontenc}
\RequirePackage[french]{babel}
\RequirePackage{etoolbox}

\newcommand{\ShorthandsOff}{\@ifundefined{shorthandoff}{}{\shorthandoff{;:!?}}}
\newcommand{\ShorthandsOn}{\@ifundefined{shorthandon}{}{\shorthandon{;:!?}}}
\ShorthandsOff
\AtBeginDocument{\ShorthandsOn}
\AtBeginEnvironment{tikzpicture}{\ShorthandsOff}

\RequirePackage{geometry}
\RequirePackage{fancyhdr}
\RequirePackage{lastpage}
\RequirePackage{multicol}
\RequirePackage{array}
\RequirePackage{enumitem}
\RequirePackage{xcolor}
\RequirePackage{needspace}

\RequirePackage{amsmath,amssymb}
\RequirePackage{mathpazo}
\DeclareMathSizes{10.95}{10.95}{8}{5.5}
\begingroup\catcode`\_=\active \catcode`\^=\active
\gdef_#1{\sb{\scriptscriptstyle #1}}
\gdef^#1{\sp{\scriptscriptstyle\lycee@moinsepais #1}}
\endgroup
\newcommand\lycee@moins{\mathbin{%
  \setbox\z@\hbox{$\m@th\scriptscriptstyle\lycee@moinsorig$}%
  \dimen@=.55\wd\z@ \advance\dimen@-.5pt
  \hbox to.75\wd\z@{\hss$\m@th\scriptscriptstyle\vcenter{\hbox{%
    \tikz\draw[line width=.5pt,line cap=round](0,0)--(\the\dimen@,0);}}$\hss}}}
\begingroup\lccode`\~=`\- \lowercase{\endgroup\let~}\lycee@moins
\newcommand\lycee@moinsepais{\mathcode`\-="8000 }
\AtBeginDocument{\mathchardef\lycee@moinsorig=\mathcode`\-\relax}
\AtBeginDocument{\catcode`\_=12 \mathcode`\_="8000
  \catcode`\^=12 \mathcode`\^="8000 }
\newcommand\lycee@exposants{%
  \fontdimen13\textfont2=.48\fontdimen6\textfont2
  \fontdimen14\textfont2=.48\fontdimen6\textfont2
  \fontdimen15\textfont2=.48\fontdimen6\textfont2
  \fontdimen13\scriptfont2=.48\fontdimen6\scriptfont2
  \fontdimen14\scriptfont2=.48\fontdimen6\scriptfont2
  \fontdimen15\scriptfont2=.48\fontdimen6\scriptfont2}
\every@math@size\expandafter{\the\every@math@size\lycee@exposants}
\newlength{\retraitcalculs}
\setlength{\retraitcalculs}{2em}

\RequirePackage{tikz}
\usetikzlibrary{arrows.meta,calc,babel,shapes.misc}
\RequirePackage[most]{tcolorbox}
\tcbuselibrary{skins,breakable}

\definecolor{coulDef}{HTML}{1F4E79}
\definecolor{coulProp}{HTML}{7030A0}
\definecolor{coulRem}{HTML}{C55A11}
\definecolor{coulEx}{HTML}{2E7D32}
\definecolor{coulExo}{HTML}{B03A2E}
\definecolor{coulPart}{HTML}{1F4E79}
\definecolor{coulMeth}{HTML}{1B6E4F}
\definecolor{coulCourbe}{HTML}{0B5ED7}
\definecolor{coulCourbeB}{HTML}{C00000}
\definecolor{coulCourbeC}{HTML}{2E7D32}
\definecolor{coulGrille}{HTML}{8C8C8C}
\definecolor{coulRep}{HTML}{1B6E4F}
\definecolor{coulBar}{HTML}{2E7D32}

\newcommand{\R}{\mathbb{R}}
\newcommand{\Rs}{\mathbb{R}^{*}}
\renewcommand{\le}{\leqslant}
\renewcommand{\ge}{\geqslant}
\renewcommand{\approx}{\simeq}
\newcommand{\vect}[1]{\overrightarrow{#1}}
\newcommand{\norme}[1]{\left\lVert #1 \right\rVert}
\newcommand{\intervalle}[2]{\left[#1\,;#2\right]}
\RequirePackage{cancel}
\renewcommand{\CancelColor}{\color{coulExo}}
\newcommand{\fois}[1]{\times{\color{coulExo}#1}}
\newcommand{\barre}[1]{\cancel{#1}}

\tikzset{grille/.style={coulGrille,line width=0.4pt}}
\newcommand{\quadrillage}[4]{\draw[grille,step=1] (#1,#2) grid (#3,#4);}
\tikzset{
  croix/.style={cross out,draw,line width=0.6pt,inner sep=0pt,minimum size=4pt},
  croixdroite/.style={croix,rotate=45,line width=0.8pt},
}
\newcommand{\pt}[3]{\path (#1) node[croix]{} node[#2,font=\small,inner sep=2.5pt]{$#3$};}

\newcommand{\lycee@repere}[4]{%
  \draw[grille,step=1] (#1,#2) grid (#3,#4);
  \draw[-{Stealth[length=2mm]},thick] (#1,0)--({#3+0.4},0) node[below left=-1pt]{$x$};
  \draw[-{Stealth[length=2mm]},thick] (0,#2)--(0,{#4+0.4}) node[below left=-1pt]{$y$};
  \node[below left,font=\scriptsize,inner sep=1.5pt] at (0,0) {$0$};}
\newcommand{\lycee@gradx}[1]{\foreach \k/\l in {#1}{%
  \draw (\k,0.07)--(\k,-0.07) node[below,font=\scriptsize,inner sep=1.5pt]{$\l$};}}
\newcommand{\lycee@grady}[1]{\foreach \k/\l in {#1}{%
  \draw (0.07,\k)--(-0.07,\k) node[left,font=\scriptsize,inner sep=1.5pt]{$\l$};}}
\newcommand{\lycee@intff}[2]{\left[#1\,;#2\right]}
\newcommand{\lycee@intfo}[2]{\left[#1\,;#2\right[}
\newcommand{\lycee@intof}[2]{\left]#1\,;#2\right]}
\newcommand{\lycee@intoo}[2]{\left]#1\,;#2\right[}
\newcommand{\lycee@ens}[1]{\left\{#1\right\}}
\AtBeginDocument{%
  \providecommand{\repere}{\lycee@repere}%
  \providecommand{\gradx}{\lycee@gradx}%
  \providecommand{\grady}{\lycee@grady}%
  \providecommand{\Cf}{\mathcal{C}\sb{\scriptscriptstyle f}}%
  \providecommand{\Cg}{\mathcal{C}\sb{\scriptscriptstyle g}}%
  \providecommand{\intff}{\lycee@intff}%
  \providecommand{\intfo}{\lycee@intfo}%
  \providecommand{\intof}{\lycee@intof}%
  \providecommand{\intoo}{\lycee@intoo}%
  \providecommand{\ens}{\lycee@ens}}

\newlist{questions}{enumerate}{3}
\setlist[questions,1]{label=\arabic*\textdegree),leftmargin=*,itemsep=2pt,topsep=3pt}
\setlist[questions,2]{label=\alph*),leftmargin=*,itemsep=1pt,topsep=2pt}
\setlist[questions,3]{label=\roman*),leftmargin=*,itemsep=1pt,topsep=1pt}
\setlist[itemize]{label=\textbullet,leftmargin=*,itemsep=2pt,topsep=3pt}

\newcommand{\besoinplace}[1]{\par\penalty\@M\begingroup
  \dimen@=#1\relax\dimen@ii=\pagegoal\advance\dimen@ii-\pagetotal
  \ifdim\dimen@>\dimen@ii\ifdim\dimen@ii>\z@\vfil\fi\break\fi\endgroup}

\newcommand{\lycee@pied}{\thepage/\pageref{LastPage}}
\newcommand{\entete}[2]{%
  \pagestyle{fancy}\fancyhf{}%
  \fancyhead[L]{\footnotesize\color{coulPart}\textbf{#1}}%
  \fancyhead[R]{\footnotesize\color{coulPart}\textbf{#2}}%
  \fancyfoot[C]{\footnotesize\lycee@pied}%
  \renewcommand{\headrulewidth}{0.4pt}%
  \renewcommand{\headrule}{\hbox to\headwidth{%
    \color{coulPart!40}\leaders\hrule height \headrulewidth\hfill}}}

\RequirePackage{ccicons}
\newcommand{\lycee@licence}{{\scriptsize\color{gray}%
  \copyright\ Luc Destombe --- \ccbyncsa\ CC BY-NC-SA 4.0}}
\newif\iflycee@licence \lycee@licencetrue
\newcommand{\sanslicence}{\lycee@licencefalse}
\AtBeginDocument{\iflycee@licence\AddToHookNext{shipout/foreground}{%
  \put(0,0){\rlap{\kern\dimexpr1in+\hoffset+\oddsidemargin+\textwidth\relax
    \llap{\raisebox{-\dimexpr1in+\voffset+\topmargin+\headheight+\headsep
      +\textheight+\footskip\relax}{\lycee@licence}}}}}\fi}

\newcommand{\formulesserrees}{%
  \setlength{\abovedisplayskip}{3pt}\setlength{\belowdisplayskip}{3pt}%
  \setlength{\abovedisplayshortskip}{2pt}\setlength{\belowdisplayshortskip}{3pt}}

\setlength{\parindent}{0pt}

\RequirePackage[implicit=false,hidelinks,pdfencoding=auto,psdextra,bookmarksopen,
  bookmarksopenlevel=1,pdfcreator={},pdfproducer={}]{hyperref}
\RequirePackage{bookmark}
\hypersetup{pdfauthor={Luc Destombe},pdfkeywords={CC BY-NC-SA 4.0}}
\newcounter{lycee@signet}
\newcommand{\signet}[3][]{\stepcounter{lycee@signet}%
  \edef\lycee@dest{\if\relax\detokenize{#1}\relax
    signet.\arabic{lycee@signet}\else#1\fi}%
  \hypertarget{\lycee@dest}{}\bookmark[level=#2,dest=\lycee@dest]{#3}}
\pdfstringdefDisableCommands{%
  \def\R{R}\def\vect#1{#1}\def\Cf{Cf}\def\Cg{Cg}\def\fois#1{×#1}%
  \def\barre#1{#1}\def\intervalle#1#2{[#1;#2]}\def\intff#1#2{[#1;#2]}%
  \def\textsuperscript#1{#1}\def\dfrac#1#2{#1/#2}\def\frac#1#2{#1/#2}%
  \def\vec#1{#1⃗}}%
\geometry{top=2cm,bottom=1cm,left=1cm,right=1cm,headheight=14pt,footskip=0.6cm}

\RequirePackage{pgfplots}
\pgfplotsset{compat=1.18}
\RequirePackage{tkz-tab}

\newcounter{partiecours}
\renewcommand{\thepartiecours}{\Roman{partiecours}}
\newcommand{\partiecours}[1]{%
  \refstepcounter{partiecours}%
  \Needspace*{8\baselineskip}%
  \bigskip
  \noindent\begin{tcolorbox}[enhanced,colback=coulPart,colframe=coulPart,
      boxrule=0pt,arc=2pt,left=8pt,right=8pt,top=4pt,bottom=4pt,
      before skip=6pt,after skip=10pt]
    \color{white}\bfseries\large \thepartiecours{} --- #1%
    \signet[partie-\arabic{partiecours}]{1}{\thepartiecours{} --- #1}
  \end{tcolorbox}%
  \addcontentsline{toc}{section}{\thepartiecours{} --- #1}%
}

\newtcolorbox{boitecours}[3][]{%
  enhanced,breakable,
  colback=white,colframe=#2,coltitle=white,
  fonttitle=\bfseries,
  attach boxed title to top left={xshift=8pt,yshift=-\tcboxedtitleheight/2},
  boxed title style={colback=#2,arc=2pt,boxrule=0pt},
  boxrule=0.9pt,arc=2pt,
  left=8pt,right=8pt,top=8pt,bottom=6pt,
  title={#3},#1}

\newcounter{methode}
\newenvironment{definitionbox}[1][Définition]%
  {\begin{boitecours}[phantom={\signet{2}{#1}}]{coulDef}{#1}}{\end{boitecours}}
\newenvironment{proprietebox}[1][Propriété]%
  {\begin{boitecours}[phantom={\signet{2}{#1}}]{coulProp}{#1}}{\end{boitecours}}
\newenvironment{methodebox}[1][Méthode]{\stepcounter{methode}%
  \begin{boitecours}[phantom={\signet[methode-\arabic{methode}]{2}{#1}}]{coulMeth}{#1}}%
  {\end{boitecours}}

\newcommand{\etape}[1]{\par\smallskip\textbf{\color{coulMeth}#1}\par\nobreak\smallskip}
\newcommand{\enonce}[1]{\emph{#1}\par\smallskip}
\newcommand{\reponse}[1]{\par\smallskip
  {\footnotesize\itshape\color{black!55}Réponses : #1}}
\newcommand{\demo}[1]{\textbf{\color{coulMeth}Démonstration.} #1}
\newcommand{\afaire}[1]{%
  \par\smallskip
  \noindent{\small\color{coulExo}$\blacktriangleright$\ \textbf{À faire :}
    fiche d'exercices du chapitre, #1.}\par\smallskip}

\newenvironment{calculs}{\setlength{\jot}{5pt}\@fleqntrue\@mathmargin\retraitcalculs
  \setlength{\abovedisplayskip}{4pt}\setlength{\belowdisplayskip}{4pt}%
  \setlength{\abovedisplayshortskip}{4pt}\setlength{\belowdisplayshortskip}{4pt}%
  \csname alignat*\endcsname{2}}%
  {\csname endalignat*\endcsname}
\newcommand{\rem}[1]{\text{\small\itshape\color{black!60}#1}}
\newcommand{\explique}[2]{\par\smallskip\noindent
  \begin{minipage}[t]{0.57\linewidth}\raggedright #1\end{minipage}\hfill
  \begin{minipage}[t]{0.40\linewidth}\raggedright\leavevmode
    {\small\itshape\color{black!60}#2}\end{minipage}\par}
\newcommand{\cas}[1]{\par\smallskip\noindent\textbf{\color{coulExo}#1}\nobreak}
\newcommand{\calc}[1]{\par\smallskip\textbf{\color{coulExo}#1}\par\nobreak}

\newtcolorbox{remarquebox}[1][Remarque]{%
  enhanced,breakable,colback=white,colframe=coulRem,
  boxrule=0pt,leftrule=2.5pt,arc=0pt,outer arc=0pt,
  left=8pt,right=6pt,top=5pt,bottom=5pt,
  before upper={\textbf{\color{coulRem}#1}\par\smallskip}}

\newtcolorbox{exemplebox}[1][Exemple]{%
  enhanced,breakable,colback=white,colframe=coulEx,
  boxrule=0pt,leftrule=2.5pt,arc=0pt,outer arc=0pt,
  left=8pt,right=6pt,top=5pt,bottom=5pt,
  before upper={\textbf{\color{coulEx}#1}\par\smallskip}}

\pgfplotsset{
  repere/.style={
    axis lines=middle,
    axis line style={-{Stealth[length=2.4mm]},thick},
    every axis x label/.style={at={(ticklabel* cs:1.0)},anchor=west},
    every axis y label/.style={at={(ticklabel* cs:1.0)},anchor=south},
    label style={font=\footnotesize},
    tick label style={font=\scriptsize},
    grid=both,
    major grid style={grille},
    minor tick num=0,
    tick align=inside,
    clip mode=individual,
  },
}
\tikzset{
  courbe/.style={coulCourbe,line width=1.1pt,smooth},
  courbeB/.style={coulCourbeB,line width=1.1pt,smooth},
  courbeC/.style={coulCourbeC,line width=1.1pt,smooth},
}

\newlist{etapes}{enumerate}{1}
\setlist[etapes]{label=\textbf{\arabic*.},leftmargin=*,itemsep=1pt,topsep=2pt}

\newlist{expressions}{enumerate}{1}
\setlist[expressions]{label={\Alph*\ $=$},leftmargin=*,itemsep=3pt,topsep=3pt}

\newcounter{exo}
\renewcommand{\theexo}{\ifnum\value{exo}<10 0\fi\arabic{exo}}
\newtcolorbox{exercice}[1][]{%
  enhanced,breakable,
  colback=white,colframe=coulExo,
  boxrule=0.9pt,arc=2pt,
  left=8pt,right=8pt,top=8pt,bottom=6pt,
  coltitle=white,fonttitle=\bfseries,
  attach boxed title to top left={xshift=8pt,yshift=-\tcboxedtitleheight/2},
  boxed title style={colback=coulExo,arc=2pt,boxrule=0pt},
  phantom={\signet{2}{Exercice \arabic{exo}}},
  title={Exercice \theexo\ifx\relax#1\relax\else{} --- #1\fi}}
\newenvironment{exo}[1][\relax]{\refstepcounter{exo}\begin{exercice}[#1]}{\end{exercice}}

  \RequirePackage{comment}
  \excludecomment{correction}
  \renewcommand{\reponse}[1]{}
  \newcommand{\autableau}{\hfill{\footnotesize\itshape\color{black!45}(au tableau)}}
  \newcommand{\etapeexemples}{%
    \par\smallskip\textbf{\color{coulMeth}Exemples}\autableau\par\nobreak\smallskip}
  \newcommand{\etapetableau}[1]{%
    \par\smallskip\textbf{\color{coulMeth}#1}\autableau\par\nobreak\smallskip}
  \newtcolorbox{exempletableau}[1][Exemple]{%
    enhanced,breakable,colback=white,colframe=coulEx,
    boxrule=0pt,leftrule=2.5pt,arc=0pt,outer arc=0pt,
    left=8pt,right=6pt,top=4pt,bottom=4pt,
    before skip=5pt,after skip=5pt,
    before upper={\textbf{\color{coulEx}#1}\autableau\par\smallskip}}

\RequirePackage{comment}
\excludecomment{correctionavj}

\newcommand{\coursEntete}{}
\newcommand{\coursNiveau}{}
\newcommand{\coursTitre}{}
\newcommand{\coursSurTitre}{}
\newcommand{\coursSousTitre}{}

\newcommand{\enteteCours}[1]{\renewcommand{\coursEntete}{#1}}
\newcommand{\chapitrecours}[2]{\enteteCours{Chapitre #1 --- #2}}
\newcommand{\niveaucours}[1]{\renewcommand{\coursNiveau}{#1}}
\newcommand{\titrecours}[3][]{%
  \renewcommand{\coursTitre}{#2}%
  \renewcommand{\coursSurTitre}{#1}%
  \renewcommand{\coursSousTitre}{#3}}

\pagestyle{fancy}
\fancyhf{}
\fancyhead[L]{\footnotesize\color{coulPart}\textbf{\coursEntete}}
\fancyhead[R]{\footnotesize\color{coulPart}\textbf{\coursNiveau}}
\fancyfoot[C]{\footnotesize\thepage}
\renewcommand{\headrulewidth}{0.4pt}
\renewcommand{\headrule}{\hbox to\headwidth{\color{coulPart!40}\leaders\hrule height \headrulewidth\hfill}}

\setlength{\parskip}{2pt}

\AtBeginDocument{%
  \begin{center}
    \begin{tcolorbox}[enhanced,width=0.72\linewidth,colback=white,colframe=coulPart,
        boxrule=1.6pt,arc=3pt,halign=center,top=10pt,bottom=10pt]
      \ifx\coursSurTitre\empty
        {\Huge\bfseries\color{coulPart} \coursTitre}\\[4pt]
      \else
        {\Huge\bfseries\color{coulPart} \coursTitre}\\[3pt]
        {\Large\bfseries\color{coulPart} \coursSurTitre}\\[5pt]
      \fi
      {\normalsize \coursSousTitre}%
      \\[2pt]{\small\itshape Version élève}
    \end{tcolorbox}
  \end{center}%
}
\makeatother

\chapitrecours{6}{Produit scalaire}
\niveaucours{1\textsuperscript{re} STI2D}
\titrecours{PRODUIT SCALAIRE}{Cours --- Première STI2D}

\begin{document}

\newcommand{\coord}[2]{\begin{pmatrix}#1\\#2\end{pmatrix}}

\ShorthandsOff
\tikzset{
  vecteur/.style={-{Stealth[length=2.6mm]},line width=1.1pt,coulCourbe},
  vecteurB/.style={-{Stealth[length=2.6mm]},line width=1.1pt,coulCourbeB},
  vecteurC/.style={-{Stealth[length=2.6mm]},line width=1.1pt,coulCourbeC},
  aide/.style={dashed,black!45,line width=0.6pt},
  fig/.style={line width=0.8pt,black!75},
  axe/.style={aide,-{Stealth[length=2.2mm]}},
  deplacement/.style={-{Stealth[length=2.6mm]},line width=1.1pt,black!75},
}
\ShorthandsOn

\newcommand{\troisforces}[6]{%
\begin{tikzpicture}[scale=0.75]
  \draw[fig,fill=black!6] (-0.7,-0.45) rectangle (0.7,0.45);
  \draw[axe] (-3.2,0) -- (6.2,0);
  \draw[deplacement] (0,0) -- node[pos=0.6,below,font=\footnotesize]{déplacement\,: #6} (5,0);
  \draw (1.5,0) arc (0:#1:1.5);
  \node[font=\footnotesize,anchor=west] at ({1.55*cos(#1/2)},{1.55*sin(#1/2)}) {#2};
  \draw[vecteur] (0,0) -- node[midway,above left,font=\footnotesize]{#3} ({3.2*cos(#1)},{3.2*sin(#1)});
  \draw[vecteurB] (0,0) -- node[pos=0.8,right,font=\footnotesize]{#4} (0,-2.4);
  \draw[vecteurC] (0,0) -- node[pos=0.8,above,font=\footnotesize]{#5} (-2.6,0);
  \pt{0,0}{below left}{G}
\end{tikzpicture}}

\newcommand{\vecteuraxe}[2]{%
\begin{tikzpicture}[scale=0.55]
  \draw[axe] (-2.6,0) -- (4.2,0) node[right,font=\footnotesize]{axe};
  \ifnum#2=1
    \draw[aide] ({3*cos(#1)},{3*sin(#1)}) -- ({3*cos(#1)},-0.35);
    \draw[red,line width=1.6pt] (0,-0.35) -- node[below,font=\footnotesize]{$p$} ({3*cos(#1)},-0.35);
  \fi
  \draw[deplacement] (0,0) -- node[pos=0.8,above,font=\footnotesize]{$\vec{u}$} (2.5,0);
  \draw (0.8,0) arc (0:#1:0.8);
  \node[font=\footnotesize,anchor=south west] at ({0.85*cos(#1/2)},{0.85*sin(#1/2)}) {$#1^{\circ}$};
  \draw[vecteur] (0,0) -- node[midway,left,font=\footnotesize]{$\vec{v}$} ({3*cos(#1)},{3*sin(#1)});
\end{tikzpicture}}

\newcommand{\solidetire}[4]{%
\begin{tikzpicture}[scale=0.75]
  \draw[fig,fill=black!6] (-0.7,-0.45) rectangle (0.7,0.45);
  \draw[axe] (-1.6,0) -- (6.2,0);
  \draw[deplacement] (0,0) -- node[pos=0.6,below,font=\footnotesize]{déplacement\,: #4} (5,0);
  \draw (1.5,0) arc (0:#1:1.5);
  \node[font=\footnotesize,anchor=west] at ({1.55*cos(#1/2)},{1.55*sin(#1/2)}) {$#1^{\circ}$};
  \draw[vecteur] (0,0) -- node[midway,above left,font=\footnotesize]{#2\,: #3} ({3.4*cos(#1)},{3.4*sin(#1)});
  \pt{0,0}{below left}{G}
\end{tikzpicture}}

\begin{remarquebox}[Comment utiliser ce cours]
Chaque partie commence par ce qu'il faut \textbf{savoir} (définitions, théorèmes), puis
vient ce qu'il faut \textbf{savoir faire} : une méthode, avec des
\textbf{exemples} faits en classe, puis des exercices
\og\textbf{À vous de jouer}\fg{} à chercher seul.

\smallskip
Sur cette feuille, seuls les \textbf{énoncés} des exemples figurent : la résolution,
signalée par la mention \emph{(au tableau)}, est à rédiger dans le cahier.

\end{remarquebox}

\partiecours{Coordonnées d'un vecteur}

\begin{definitionbox}[Définition 1 --- Coordonnées d'un vecteur]
\begin{minipage}[t]{0.58\linewidth}\raggedright
\begin{itemize}
  \item Dans un repère, un vecteur décrit un \textbf{déplacement}.
  \item Les nombres $x$ et $y$ sont les \textbf{coordonnées} du vecteur. On les écrit
        en colonne, $\vec{u}\coord{x}{y}$, ou en ligne, $\vec{u}\,(x\,;y)$.
  \item Sur la figure, pour aller de $A$ à $B$, on avance de $3$ vers la droite et de
        $2$ vers le haut : $\vect{AB}\coord{3}{2}$.
\end{itemize}
\end{minipage}\hfill
\begin{minipage}[t]{0.38\linewidth}
\vspace{0pt}\centering
\begin{tikzpicture}[x=0.6cm,y=0.6cm]
  \quadrillage{0.4}{0.4}{4.6}{3.6}
  \draw[aide] (1,1) -- (4,1) -- (4,3);
  \draw[vecteur] (1,1) -- node[midway,above left,font=\footnotesize]{$\vect{AB}$} (4,3);
  \node[below,font=\footnotesize] at (2.5,1) {$3$};
  \node[right,font=\footnotesize] at (4,2) {$2$};
  \pt{1,1}{below left}{A} \pt{4,3}{above right}{B}
\end{tikzpicture}
\end{minipage}
\end{definitionbox}

\begin{proprietebox}[Propriété 1 --- Calcul des coordonnées]
Pour deux points $A(x_{A}\,;y_{A})$ et $B(x_{B}\,;y_{B})$, le vecteur $\vect{AB}$ a pour
coordonnées
\[
  \vect{AB}\coord{x_{B}-x_{A}}{y_{B}-y_{A}}
  \qquad\text{(arrivée moins départ).}
\]
\end{proprietebox}

\begin{methodebox}[Méthode 1 --- Lire et tracer un vecteur dans un repère]
\etapeexemples
\begin{minipage}[t]{0.58\linewidth}\raggedright
Dans le repère ci-contre :
\begin{questions}[label=\alph*)]
  \item lire les coordonnées des vecteurs $\vect{AB}$ et $\vect{CD}$ ;
  \item tracer un représentant du vecteur $\vec{u}\coord{3}{-1}$ d'origine $E$.
\end{questions}
\end{minipage}\hfill
\begin{minipage}[t]{0.38\linewidth}
\vspace{0pt}\centering
\begin{tikzpicture}[x=0.45cm,y=0.45cm]
  \repere{-4}{-4}{4}{4}
  \gradx{-3,-2,-1,1,2,3} \grady{-3,-2,-1,1,2,3}
  \pt{-3,1}{above left}{A} \pt{1,3}{above}{B} \pt{3,0}{above}{C} \pt{0,-2}{right}{D}
  \pt{-3,-2}{left}{E}
\end{tikzpicture}
\end{minipage}

\etape{À vous de jouer}
\begin{minipage}[t]{0.58\linewidth}\raggedright
Dans le repère ci-contre :
\begin{questions}[label=\alph*)]
  \item lire les coordonnées des vecteurs $\vect{FG}$ et $\vect{HK}$ ;
  \item tracer un représentant du vecteur $\vec{w}\coord{-2}{3}$ d'origine $L$.
\end{questions}
\end{minipage}\hfill
\begin{minipage}[t]{0.38\linewidth}
\vspace{0pt}\centering
\begin{tikzpicture}[x=0.45cm,y=0.45cm]
  \repere{-4}{-4}{4}{4}
  \gradx{-3,-2,-1,1,2,3} \grady{-3,-2,-1,1,2,3}
  \pt{-2,3}{above left}{F} \pt{2,2}{above right}{G} \pt{-3,-1}{below left}{H} \pt{0,-3}{right}{K}
  \pt{3,-3}{right}{L}
\end{tikzpicture}
\end{minipage}

\end{methodebox}

\begin{methodebox}[Méthode 2 --- Calculer les coordonnées d'un vecteur]
\etapeexemples
On donne les points $A(2\,;-1)$, $B(5\,;3)$, $C(-4\,;2)$ et $D(1\,;-3)$. Calculer les
coordonnées des vecteurs :
\begin{questions}[label=\alph*)]
  \item $\vect{AB}$ ;
  \item $\vect{CD}$ ;
  \item $\vect{BA}$.
\end{questions}

\etape{À vous de jouer}
On donne les points $E(1\,;4)$, $F(3\,;-2)$, $G(-3\,;-1)$ et $H(2\,;5)$. Calculer les
coordonnées des vecteurs :
\begin{questions}[label=\alph*)]
  \item $\vect{EF}$ ;
  \item $\vect{GH}$ ;
  \item $\vect{FE}$.
\end{questions}

\end{methodebox}

\afaire{exercices 1 à 4}

\partiecours{Somme de vecteurs, produit par un nombre réel}

\begin{proprietebox}[Propriété 2 --- Calculs sur les coordonnées]
Pour deux vecteurs $\vec{u}\coord{x}{y}$ et $\vec{v}\coord{x'}{y'}$ et un nombre réel
$k$ :
\[
  \vec{u}+\vec{v}\coord{x+x'}{y+y'}
  \qquad
  \vec{u}-\vec{v}\coord{x-x'}{y-y'}
  \qquad
  k\vec{u}\coord{kx}{ky}.
\]
On ajoute (ou on multiplie) les abscisses entre elles et les ordonnées entre elles.
\end{proprietebox}

\begin{remarquebox}[Un exemple concret : deux forces]
Deux câbles tirent une charge avec les forces $\vec{F_{1}}\coord{3}{1}$ et
$\vec{F_{2}}\coord{-1}{2}$ (en kN). L'effet est le même qu'une force unique, la
\textbf{résultante} $\vec{F_{1}}+\vec{F_{2}}\coord{2}{3}$.
\end{remarquebox}

\begin{methodebox}[Méthode 3 --- Calculer avec les coordonnées]
\etapeexemples
On donne les vecteurs $\vec{u}\coord{3}{-2}$ et $\vec{v}\coord{-1}{4}$, et les points
$A(1\,;2)$, $B(4\,;0)$ et $C(-1\,;5)$.
\begin{questions}[label=\alph*)]
  \item Calculer les coordonnées du vecteur $\vec{u}+\vec{v}$.
  \item Calculer les coordonnées du vecteur $3\vec{u}-2\vec{v}$.
  \item Calculer les coordonnées du vecteur $\vect{AB}+2\vect{AC}$.
  \item Calculer les coordonnées du vecteur $\vec{w}$ tel que $\vec{u}+\vec{w}\coord{5}{1}$.
\end{questions}

\etape{À vous de jouer}
On donne les vecteurs $\vec{u}\coord{-2}{5}$ et $\vec{v}\coord{4}{1}$, et les points
$E(0\,;3)$, $F(2\,;-1)$ et $G(3\,;1)$.
\begin{questions}[label=\alph*)]
  \item Calculer les coordonnées du vecteur $\vec{u}+\vec{v}$.
  \item Calculer les coordonnées du vecteur $3\vec{u}-2\vec{v}$.
  \item Calculer les coordonnées du vecteur $\vect{EF}+2\vect{EG}$.
  \item Calculer les coordonnées du vecteur $\vec{w}$ tel que $\vec{u}+\vec{w}\coord{1}{2}$.
\end{questions}

\end{methodebox}

\afaire{exercices 5 à 8}

\partiecours{Parallélogramme et coordonnées d'un point}

\begin{proprietebox}[Propriété 3 --- Égalité de vecteurs et parallélogramme]
\begin{minipage}[t]{0.58\linewidth}\raggedright
Deux vecteurs sont \textbf{égaux} quand ils ont les mêmes coordonnées : même
abscisse et même ordonnée.

\smallskip
$ABCD$ est un \textbf{parallélogramme} quand $\vect{AB}=\vect{DC}$ (les côtés
$[AB]$ et $[DC]$ sont parallèles, de même longueur et parcourus dans le même sens).
\end{minipage}\hfill
\begin{minipage}[t]{0.38\linewidth}
\vspace{0pt}\centering
\begin{tikzpicture}[scale=0.55]
  \draw[fig] (0,0) -- (4,1) -- (5,3) -- (1,2) -- cycle;
  \draw[vecteur] (0,0) -- node[midway,below right,font=\footnotesize]{$\vect{AB}$} (4,1);
  \draw[vecteurB] (1,2) -- node[midway,above left,font=\footnotesize]{$\vect{DC}$} (5,3);
  \pt{0,0}{below left}{A} \pt{4,1}{below right}{B} \pt{5,3}{above right}{C}
  \pt{1,2}{above left}{D}
\end{tikzpicture}
\end{minipage}
\end{proprietebox}

\begin{methodebox}[Méthode 4 --- Calculer les coordonnées d'un point défini par une égalité de vecteurs]
L'égalité de deux vecteurs donne deux équations : une pour les abscisses, une pour
les ordonnées.
\etapeexemples
On donne les points $A(1\,;2)$, $B(5\,;3)$, $C(4\,;-1)$ et le vecteur
$\vec{u}\coord{3}{-1}$.
\begin{questions}[label=\alph*)]
  \item Calculer les coordonnées du point $E$ tel que $\vect{AE}=\vec{u}$.
  \item Calculer les coordonnées du point $D$ tel que $ABCD$ soit un parallélogramme.
\end{questions}

\etape{À vous de jouer}
On donne les points $E(-2\,;1)$, $F(2\,;4)$, $G(3\,;-1)$ et le vecteur
$\vec{w}\coord{2}{-3}$.
\begin{questions}[label=\alph*)]
  \item Calculer les coordonnées du point $H$ tel que $\vect{EH}=\vec{w}$.
  \item Calculer les coordonnées du point $K$ tel que $EFGK$ soit un parallélogramme.
\end{questions}

\end{methodebox}

\afaire{exercices 9 à 11}

\partiecours{Colinéarité de deux vecteurs}

\begin{definitionbox}[Définition 2 --- Vecteurs colinéaires]
Deux vecteurs $\vec{u}$ et $\vec{v}$ sont \textbf{colinéaires} quand l'un est un multiple
de l'autre : $\vec{v}=k\vec{u}$ pour un nombre réel $k$. Ils ont alors la même
direction : de \textbf{même sens} si $k>0$, de \textbf{sens contraires} si $k<0$.

\smallskip
Exemple : $\vec{u}\coord{2}{1}$ et $\vec{v}\coord{6}{3}$ sont colinéaires et de même
sens, car $\vec{v}=3\vec{u}$ ; $\vec{w}\coord{-4}{-2}$ est colinéaire à $\vec{u}$, de sens
contraire, car $\vec{w}=-2\vec{u}$.
\end{definitionbox}

\begin{proprietebox}[Propriété 4 --- Critère de colinéarité]
Pour deux vecteurs $\vec{u}\coord{x}{y}$ et $\vec{v}\coord{x'}{y'}$, le nombre
\[
  \det(\vec{u}\,;\vec{v})=xy'-yx'
\]
s'appelle le \textbf{déterminant} de $\vec{u}$ et $\vec{v}$ (produit en croix des
coordonnées).

\smallskip
$\vec{u}$ et $\vec{v}$ sont colinéaires quand $\det(\vec{u}\,;\vec{v})=0$ ; ils ne le
sont pas quand $\det(\vec{u}\,;\vec{v})\neq 0$.
\end{proprietebox}

\begin{methodebox}[Méthode 5 --- Démontrer que deux vecteurs sont colinéaires, ou non]
On cherche d'abord un multiple évident ($\vec{v}=k\vec{u}$, définition 2) ; s'il ne se
voit pas, on calcule le déterminant (propriété 4).
\etapeexemples
Les vecteurs suivants sont-ils colinéaires ?
\begin{questions}[label=\alph*)]
  \item $\vec{u}\coord{2}{3}$ et $\vec{v}\coord{6}{9}$ ;
  \item $\vec{u}\coord{4}{6}$ et $\vec{w}\coord{6}{10}$ ;
  \item $\vec{u}\coord{1}{-2}$ et $\vec{r}\coord{-4}{8}$ ;
  \item $\vec{u}\coord{4}{6}$ et $\vec{s}\coord{6}{9}$ ;
  \item Calculer le nombre réel $m$ pour que les vecteurs $\vec{u}\coord{2}{5}$ et
        $\vec{t}\coord{6}{m}$ soient colinéaires.
\end{questions}

\etape{À vous de jouer}
Les vecteurs suivants sont-ils colinéaires ?
\begin{questions}[label=\alph*)]
  \item $\vec{u}\coord{3}{1}$ et $\vec{v}\coord{12}{4}$ ;
  \item $\vec{u}\coord{3}{7}$ et $\vec{w}\coord{5}{12}$ ;
  \item $\vec{u}\coord{-2}{5}$ et $\vec{r}\coord{6}{-15}$ ;
  \item $\vec{u}\coord{6}{4}$ et $\vec{s}\coord{9}{6}$ ;
  \item Calculer le nombre réel $m$ pour que les vecteurs $\vec{u}\coord{3}{-2}$ et
        $\vec{t}\coord{9}{m}$ soient colinéaires.
\end{questions}

\end{methodebox}

\afaire{exercices 12 à 15}

\partiecours{Droites parallèles, points alignés}

\begin{proprietebox}[Propriété 5 --- Parallélisme et alignement]
\begin{itemize}
  \item Les droites $(AB)$ et $(CD)$ sont \textbf{parallèles} quand les vecteurs
        $\vect{AB}$ et $\vect{CD}$ sont colinéaires.
  \item Les points $A$, $B$ et $C$ sont \textbf{alignés} quand les vecteurs
        $\vect{AB}$ et $\vect{AC}$ sont colinéaires.
\end{itemize}
\end{proprietebox}

\begin{methodebox}[Méthode 6 --- Démontrer un parallélisme ou un alignement]
Une droite donnée par une équation cartésienne $ax+by+c=0$ (vue en début d'année) a pour
vecteur directeur $\vec{u}\coord{-b}{a}$ : deux telles droites sont parallèles quand
leurs vecteurs directeurs sont colinéaires.
\etapeexemples
On donne les points $A(-2\,;1)$, $B(2\,;3)$, $C(0\,;-2)$, $D(6\,;1)$ et $E(6\,;5)$.
\begin{questions}[label=\alph*)]
  \item Les droites $(AB)$ et $(CD)$ sont-elles parallèles ?
  \item Les points $A$, $B$ et $E$ sont-ils alignés ?
  \item Les points $A$, $B$ et $C$ sont-ils alignés ?
  \item Les droites $(d_{1})$ : $2x-3y+1=0$ et $(d_{2})$ : $-4x+6y+5=0$ sont-elles
        parallèles ?
\end{questions}

\etape{À vous de jouer}
On donne les points $F(-3\,;2)$, $G(1\,;0)$, $H(-1\,;4)$, $K(5\,;1)$ et $L(5\,;-2)$.
\begin{questions}[label=\alph*)]
  \item Les droites $(FG)$ et $(HK)$ sont-elles parallèles ?
  \item Les points $F$, $G$ et $L$ sont-ils alignés ?
  \item Les points $F$, $G$ et $H$ sont-ils alignés ?
  \item Les droites $(d_{3})$ : $x+2y-3=0$ et $(d_{4})$ : $3x+6y+2=0$ sont-elles
        parallèles ?
\end{questions}

\end{methodebox}

\afaire{exercices 16 à 19}

\partiecours{Coordonnées du milieu d'un segment}

\begin{proprietebox}[Propriété 6 --- Milieu]
Le milieu $I$ du segment $[AB]$ a pour coordonnées les \textbf{moyennes} des
coordonnées de $A$ et de $B$ :
\[
  x_{I}=\frac{x_{A}+x_{B}}{2}
  \qquad\text{et}\qquad
  y_{I}=\frac{y_{A}+y_{B}}{2}.
\]
\end{proprietebox}

\begin{methodebox}[Méthode 7 --- Utiliser les coordonnées du milieu]
\etapeexemples
On donne les points $A(-1\,;4)$, $B(5\,;2)$ et $K(1\,;1)$.
\begin{questions}[label=\alph*)]
  \item Calculer les coordonnées du milieu $I$ de $[AB]$.
  \item $K$ est le milieu du segment $[AC]$. Calculer les coordonnées du point $C$.
\end{questions}

\etape{À vous de jouer}
On donne les points $E(3\,;-2)$, $F(-5\,;6)$ et $M(2\,;1)$.
\begin{questions}[label=\alph*)]
  \item Calculer les coordonnées du milieu $J$ de $[EF]$.
  \item $M$ est le milieu du segment $[EG]$. Calculer les coordonnées du point $G$.
\end{questions}

\end{methodebox}

\afaire{exercices 20 à 22}

\partiecours{Norme d'un vecteur}

\begin{definitionbox}[Définition 3 --- Norme]
\begin{minipage}[t]{0.58\linewidth}\raggedright
La \textbf{norme} d'un vecteur $\vec{u}$, notée $\norme{\vec{u}}$, est sa
\textbf{longueur}. Pour deux points $A$ et $B$ : $\norme{\vect{AB}}=AB$.

\smallskip
Dans un repère \textbf{orthonormé}, pour $\vec{u}\coord{x}{y}$, le théorème de
Pythagore donne
\[
  \norme{\vec{u}}=\sqrt{x^{2}+y^{2}}.
\]
\end{minipage}\hfill
\begin{minipage}[t]{0.38\linewidth}
\vspace{0pt}\centering
\begin{tikzpicture}[x=0.6cm,y=0.6cm]
  \quadrillage{-0.4}{-0.4}{4.4}{3.4}
  \draw[aide] (0,0) -- (4,0) -- (4,3);
  \draw[black!60] (3.7,0) -- (3.7,0.3) -- (4,0.3);
  \draw[vecteur] (0,0) -- node[midway,above left,font=\footnotesize]{$\vec{u}$} (4,3);
  \node[below,font=\footnotesize] at (2,0) {$x$};
  \node[right,font=\footnotesize] at (4,1.5) {$y$};
\end{tikzpicture}
\end{minipage}
\end{definitionbox}

\begin{methodebox}[Méthode 8 --- Calculer une norme, une distance]
Le repère est orthonormé.
\etapeexemples
\begin{questions}[label=\alph*)]
  \item Calculer la norme du vecteur $\vec{u}\coord{3}{-4}$.
  \item On donne $A(2\,;1)$ et $B(4\,;5)$. Calculer la distance $AB$, puis en donner
        une valeur approchée au centième.
  \item On donne $A(1\,;1)$, $B(5\,;3)$ et $C(3\,;5)$. Quelle est la nature du
        triangle $ABC$ ?
\end{questions}

\etape{À vous de jouer}
Le repère est orthonormé.
\begin{questions}[label=\alph*)]
  \item Calculer la norme du vecteur $\vec{v}\coord{-6}{8}$.
  \item On donne $E(-2\,;1)$ et $F(1\,;3)$. Calculer la distance $EF$, puis en donner
        une valeur approchée au centième.
  \item On donne $K(-1\,;2)$, $L(3\,;4)$ et $M(1\,;-2)$. Quelle est la nature du
        triangle $KLM$ ?
\end{questions}

\end{methodebox}

\afaire{exercices 23 à 25}

\partiecours{Définition du produit scalaire}

\begin{remarquebox}[Un problème concret : le travail d'une force]
\begin{minipage}[t]{0.58\linewidth}\raggedright
Un chariot est tiré sur $10$~m par une corde qui exerce une force $\vec{F}$ de $50$~N,
inclinée de $30^{\circ}$ par rapport au sol. Seule la partie de la force dirigée le long
du sol fait avancer le chariot. En physique, le travail de la force vaut
\[
  W=50\times 10\times\cos(30^{\circ})\approx 433\text{ J}.
\]
Ce calcul, « longueur $\times$ longueur $\times$ cosinus de l'angle », est un
\textbf{produit scalaire}.
\end{minipage}\hfill
\begin{minipage}[t]{0.38\linewidth}
\vspace{0pt}\centering
\solidetire{30}{$\vec{F}$}{$50$~N}{$10$~m}
\end{minipage}
\end{remarquebox}

\begin{definitionbox}[Définition 4 --- Produit scalaire]
\begin{minipage}[t]{0.58\linewidth}\raggedright
Le \textbf{produit scalaire} des vecteurs $\vect{AB}$ et $\vect{AC}$, de même origine
$A$, est le \textbf{nombre}
\[
  \vect{AB}\cdot\vect{AC}=AB\times AC\times\cos\big(\widehat{BAC}\big).
\]
Pour deux vecteurs $\vec{u}$ et $\vec{v}$, qui forment un angle $\alpha$ :
$\vec{u}\cdot\vec{v}=\norme{\vec{u}}\times\norme{\vec{v}}\times\cos(\alpha)$.
Si l'un des deux vecteurs est nul, $\vec{u}\cdot\vec{v}=0$.
\end{minipage}\hfill
\begin{minipage}[t]{0.38\linewidth}
\vspace{0pt}\centering
\begin{tikzpicture}[scale=0.8]
  \draw (0.7,0) arc (0:38.66:0.7);
  \node[font=\footnotesize] at (1.05,0.35) {$\alpha$};
  \draw[vecteur] (0,0) -- node[midway,below,font=\footnotesize]{$\vect{AB}$} (3.5,0);
  \draw[vecteurB] (0,0) -- node[midway,above left,font=\footnotesize]{$\vect{AC}$} (2.5,2);
  \pt{0,0}{left}{A} \pt{3.5,0}{below right}{B} \pt{2.5,2}{above right}{C}
\end{tikzpicture}
\end{minipage}

\smallskip
Le cosinus se calcule à la calculatrice, en mode \textbf{degré} :
$\cos(60^{\circ})=0{,}5$ et $\cos(120^{\circ})=-0{,}5$.

\smallskip
\textbf{Carré scalaire} : $\vec{u}\cdot\vec{u}=\norme{\vec{u}}^{2}$ (l'angle est nul et
$\cos(0^{\circ})=1$) ; on note $\vec{u}^{2}=\vec{u}\cdot\vec{u}$. Par exemple,
$\vect{AB}\cdot\vect{AB}=AB^{2}$.
\end{definitionbox}

\begin{proprietebox}[Propriété 7 --- Orthogonalité]
Deux vecteurs non nuls sont \textbf{orthogonaux} quand leur produit scalaire est
\textbf{nul} : $\vec{u}\cdot\vec{v}=0$ (car $\cos(90^{\circ})=0$).
\end{proprietebox}

\begin{methodebox}[Méthode 9 --- Calculer un produit scalaire avec un angle]
\etapeexemples
\begin{questions}[label=\alph*)]
  \item $\norme{\vec{u}}=4$, $\norme{\vec{v}}=5$ et l'angle entre $\vec{u}$ et $\vec{v}$
        mesure $60^{\circ}$. Calculer $\vec{u}\cdot\vec{v}$.
  \item Dans le triangle $ABC$, $AB=6$, $AC=3$ et $\widehat{BAC}=120^{\circ}$. Calculer
        $\vect{AB}\cdot\vect{AC}$.
  \item $EFG$ est un triangle équilatéral de côté $4$. Calculer $\vect{EF}\cdot\vect{EG}$.
\end{questions}
\begin{minipage}[t]{0.58\linewidth}\raggedright
\begin{questions}[label=\alph*),start=4]
  \item Un treuil tire une charge sur $15$~m avec une force de $200$~N qui fait un angle
        de $25^{\circ}$ avec le déplacement (schéma ci-contre). Calculer le travail de la
        force, au dixième.
\end{questions}
\end{minipage}\hfill
\begin{minipage}[t]{0.38\linewidth}
\vspace{0pt}\centering
\solidetire{25}{$\vec{F}$}{$200$~N}{$15$~m}
\end{minipage}

\etape{À vous de jouer}
\begin{questions}[label=\alph*)]
  \item $\norme{\vec{u}}=6$, $\norme{\vec{v}}=3$ et l'angle entre $\vec{u}$ et $\vec{v}$
        mesure $60^{\circ}$. Calculer $\vec{u}\cdot\vec{v}$.
  \item Dans le triangle $RST$, $RS=8$, $RT=5$ et $\widehat{SRT}=120^{\circ}$. Calculer
        $\vect{RS}\cdot\vect{RT}$.
  \item $KLM$ est un triangle équilatéral de côté $6$. Calculer $\vect{KL}\cdot\vect{KM}$.
\end{questions}
\begin{minipage}[t]{0.58\linewidth}\raggedright
\begin{questions}[label=\alph*),start=4]
  \item Un câble tire un wagonnet sur $40$~m avec une force de $150$~N qui fait un angle
        de $30^{\circ}$ avec le déplacement (schéma ci-contre). Calculer le travail de la
        force, au dixième.
\end{questions}
\end{minipage}\hfill
\begin{minipage}[t]{0.38\linewidth}
\vspace{0pt}\centering
\solidetire{30}{$\vec{F}$}{$150$~N}{$40$~m}
\end{minipage}

\end{methodebox}

\afaire{exercices 26 à 28}

\besoinplace{14\baselineskip}
\partiecours{Produit scalaire et projeté orthogonal}

\begin{definitionbox}[Définition 5 --- Projeté orthogonal]
\begin{minipage}[t]{0.58\linewidth}\raggedright
Le \textbf{projeté orthogonal} d'un point $M$ sur une droite $(d)$ est le point $H$ où
la perpendiculaire à $(d)$ passant par $M$ coupe $(d)$ : c'est le « pied » de $M$ sur
$(d)$, comme l'ombre de $M$ quand le soleil est à la verticale.
\end{minipage}\hfill
\begin{minipage}[t]{0.38\linewidth}
\vspace{0pt}\centering
\begin{tikzpicture}[scale=0.7]
  \draw[fig] (-0.5,0) -- (4.5,0) node[right,font=\footnotesize]{$(d)$};
  \draw[aide] (2,2) -- (2,0);
  \draw[black!60] (2,0.25) -- (2.25,0.25) -- (2.25,0);
  \pt{2,2}{above}{M} \pt{2,0}{below}{H}
\end{tikzpicture}
\end{minipage}
\end{definitionbox}

\begin{proprietebox}[Propriété 8 --- Produit scalaire par projection]
\begin{minipage}[t]{0.58\linewidth}\raggedright
$H$ est le projeté orthogonal de $C$ sur la droite $(AB)$. Alors
\[
  \vect{AB}\cdot\vect{AC}=\vect{AB}\cdot\vect{AH},
\]
c'est-à-dire :
\begin{itemize}
  \item $\vect{AB}\cdot\vect{AC}=AB\times AH$ si $\vect{AB}$ et $\vect{AH}$ ont le
        \textbf{même sens} ;
  \item $\vect{AB}\cdot\vect{AC}=-AB\times AH$ s'ils sont de \textbf{sens contraires}.
\end{itemize}
\end{minipage}\hfill
\begin{minipage}[t]{0.38\linewidth}
\vspace{0pt}\centering
\begin{tikzpicture}[scale=0.7]
  \draw[aide] (2.5,2) -- (2.5,0);
  \draw[black!60] (2.5,0.25) -- (2.25,0.25) -- (2.25,0);
  \draw[vecteur] (0,0) -- node[midway,below,font=\footnotesize]{$\vect{AB}$} (4,0);
  \draw[vecteurB] (0,0) -- node[midway,above left,font=\footnotesize]{$\vect{AC}$} (2.5,2);
  \pt{0,0}{left}{A} \pt{4,0}{right}{B} \pt{2.5,2}{above}{C} \pt{2.5,0}{below}{H}
\end{tikzpicture}
\end{minipage}

\smallskip
\textbf{Projection d'un vecteur sur un axe.} La projection de $\vec{v}$ sur un axe dirigé
par $\vec{u}$ est le nombre $\norme{\vec{v}}\times\cos(\alpha)$, où $\alpha$ est l'angle
entre $\vec{u}$ et $\vec{v}$ : c'est la longueur de l'« ombre » de $\vec{v}$ sur l'axe,
comptée négativement si l'angle est obtus. Alors
$\vec{u}\cdot\vec{v}=\norme{\vec{u}}\times\big(\norme{\vec{v}}\times\cos(\alpha)\big)$ :
le travail d'une force est le déplacement multiplié par la projection de la force.

\smallskip
\demo{(angle aigu) dans le triangle $AHC$ rectangle en $H$ :}
\begin{calculs}
\cos\big(\widehat{BAC}\big) &= \frac{AH}{AC} &\qquad& \rem{côté adjacent sur hypoténuse}\\
\vect{AB}\cdot\vect{AC} &= AB\times AC\times\frac{AH}{AC} && \rem{définition 4}\\
\vect{AB}\cdot\vect{AC} &= AB\times AH && \rem{on simplifie par $AC$}
\end{calculs}
\end{proprietebox}

\begin{methodebox}[Méthode 10 --- Calculer un produit scalaire par projection]
On repère le projeté orthogonal de l'extrémité du second vecteur sur la droite qui
porte le premier.
\etapeexemples
\begin{minipage}[t]{0.58\linewidth}\raggedright
$ABCD$ est un rectangle avec $AB=6$ et $AD=4$. Calculer :
\begin{questions}[label=\alph*)]
  \item $\vect{AB}\cdot\vect{AC}$ ;
  \item $\vect{AD}\cdot\vect{AC}$ ;
  \item $\vect{AB}\cdot\vect{AD}$.
\end{questions}
\end{minipage}\hfill
\begin{minipage}[t]{0.38\linewidth}
\vspace{0pt}\centering
\begin{tikzpicture}[scale=0.4]
  \draw[fig] (0,0) -- (6,0) -- (6,4) -- (0,4) -- cycle;
  \draw[aide] (0,0) -- (6,4);
  \pt{0,0}{below left}{A} \pt{6,0}{below right}{B} \pt{6,4}{above right}{C}
  \pt{0,4}{above left}{D}
\end{tikzpicture}
\end{minipage}

\smallskip
\begin{minipage}[t]{0.58\linewidth}\raggedright
\begin{questions}[label=\alph*),start=4]
  \item Dans le triangle $EFG$ ci-contre, $H$ est le projeté orthogonal de $G$ sur la
        droite $(EF)$, avec $EF=5$ et $EH=2$. Calculer $\vect{EF}\cdot\vect{EG}$.
\end{questions}
\end{minipage}\hfill
\begin{minipage}[t]{0.38\linewidth}
\vspace{0pt}\centering
\begin{tikzpicture}[scale=0.45]
  \draw[aide] (-2,0) -- (0,0);
  \draw[aide] (-2,2.5) -- (-2,0);
  \draw[black!60] (-2,0.35) -- (-1.65,0.35) -- (-1.65,0);
  \draw[fig] (0,0) -- (5,0) -- (-2,2.5) -- cycle;
  \pt{0,0}{below}{E} \pt{5,0}{below right}{F} \pt{-2,2.5}{above}{G} \pt{-2,0}{below}{H}
\end{tikzpicture}
\end{minipage}

\smallskip
\begin{minipage}[t]{0.58\linewidth}\raggedright
\begin{questions}[label=\alph*),start=5]
  \item Sur un axe dirigé par $\vec{u}$, avec $\norme{\vec{u}}=5$, le vecteur $\vec{v}$
        de norme $6$ fait un angle de $60^{\circ}$ avec $\vec{u}$ (schéma ci-contre).
        Calculer la projection $p$ de $\vec{v}$ sur l'axe, puis $\vec{u}\cdot\vec{v}$.
\end{questions}
\end{minipage}\hfill
\begin{minipage}[t]{0.38\linewidth}
\vspace{0pt}\centering
\vecteuraxe{60}{0}
\end{minipage}

\smallskip
\begin{minipage}[t]{0.58\linewidth}\raggedright
\begin{questions}[label=\alph*),start=6]
  \item Même question avec $\norme{\vec{u}}=2$, $\norme{\vec{v}}=4$ et un angle de
        $120^{\circ}$ (schéma ci-contre).
\end{questions}
\end{minipage}\hfill
\begin{minipage}[t]{0.38\linewidth}
\vspace{0pt}\centering
\vecteuraxe{120}{0}
\end{minipage}

\etape{À vous de jouer}
\begin{minipage}[t]{0.58\linewidth}\raggedright
$KLMN$ est un rectangle avec $KL=8$ et $KN=3$. Calculer :
\begin{questions}[label=\alph*)]
  \item $\vect{KL}\cdot\vect{KM}$ ;
  \item $\vect{KN}\cdot\vect{KM}$ ;
  \item $\vect{KL}\cdot\vect{KN}$.
\end{questions}
\end{minipage}\hfill
\begin{minipage}[t]{0.38\linewidth}
\vspace{0pt}\centering
\begin{tikzpicture}[scale=0.35]
  \draw[fig] (0,0) -- (8,0) -- (8,3) -- (0,3) -- cycle;
  \draw[aide] (0,0) -- (8,3);
  \pt{0,0}{below left}{K} \pt{8,0}{below right}{L} \pt{8,3}{above right}{M}
  \pt{0,3}{above left}{N}
\end{tikzpicture}
\end{minipage}

\smallskip
\begin{minipage}[t]{0.58\linewidth}\raggedright
\begin{questions}[label=\alph*),start=4]
  \item Dans le triangle $RST$ ci-contre, $H$ est le projeté orthogonal de $T$ sur la
        droite $(RS)$, avec $RS=4$ et $RH=3$. Calculer $\vect{RS}\cdot\vect{RT}$.
\end{questions}
\end{minipage}\hfill
\begin{minipage}[t]{0.38\linewidth}
\vspace{0pt}\centering
\begin{tikzpicture}[scale=0.45]
  \draw[aide] (-3,0) -- (0,0);
  \draw[aide] (-3,2) -- (-3,0);
  \draw[black!60] (-3,0.35) -- (-2.65,0.35) -- (-2.65,0);
  \draw[fig] (0,0) -- (4,0) -- (-3,2) -- cycle;
  \pt{0,0}{below}{R} \pt{4,0}{below right}{S} \pt{-3,2}{above}{T} \pt{-3,0}{below}{H}
\end{tikzpicture}
\end{minipage}

\smallskip
\begin{minipage}[t]{0.58\linewidth}\raggedright
\begin{questions}[label=\alph*),start=5]
  \item Sur un axe dirigé par $\vec{u}$, avec $\norme{\vec{u}}=3$, le vecteur $\vec{v}$
        de norme $8$ fait un angle de $60^{\circ}$ avec $\vec{u}$ (schéma ci-contre).
        Calculer la projection $p$ de $\vec{v}$ sur l'axe, puis $\vec{u}\cdot\vec{v}$.
\end{questions}
\end{minipage}\hfill
\begin{minipage}[t]{0.38\linewidth}
\vspace{0pt}\centering
\vecteuraxe{60}{0}
\end{minipage}

\smallskip
\begin{minipage}[t]{0.58\linewidth}\raggedright
\begin{questions}[label=\alph*),start=6]
  \item Même question avec $\norme{\vec{u}}=5$, $\norme{\vec{v}}=6$ et un angle de
        $120^{\circ}$ (schéma ci-contre).
\end{questions}
\end{minipage}\hfill
\begin{minipage}[t]{0.38\linewidth}
\vspace{0pt}\centering
\vecteuraxe{120}{0}
\end{minipage}

\end{methodebox}

\afaire{exercices 29 à 32}

\partiecours{Produit scalaire dans un repère orthonormé}

\begin{proprietebox}[Propriété 9 --- Expression avec les coordonnées]
Dans un repère \textbf{orthonormé}, pour $\vec{u}\coord{x}{y}$ et $\vec{v}\coord{x'}{y'}$ :
\[
  \vec{u}\cdot\vec{v}=xx'+yy'
  \qquad\text{(abscisses multipliées, plus ordonnées multipliées).}
\]
\end{proprietebox}

\begin{methodebox}[Méthode 11 --- Calculer un produit scalaire avec les coordonnées]
\etapeexemples
Le repère est orthonormé.
\begin{questions}[label=\alph*)]
  \item Calculer $\vec{u}\cdot\vec{v}$ pour $\vec{u}\coord{3}{-2}$ et $\vec{v}\coord{4}{5}$.
  \item Calculer $\vec{a}\cdot\vec{b}$ pour $\vec{a}\coord{-1}{3}$ et $\vec{b}\coord{6}{2}$.
        Que peut-on dire de ces vecteurs ?
  \item On donne $A(1\,;2)$, $B(4\,;-1)$ et $C(3\,;5)$. Calculer $\vect{AB}\cdot\vect{AC}$.
\end{questions}

\etape{À vous de jouer}
Le repère est orthonormé.
\begin{questions}[label=\alph*)]
  \item Calculer $\vec{u}\cdot\vec{v}$ pour $\vec{u}\coord{5}{-1}$ et $\vec{v}\coord{2}{3}$.
  \item Calculer $\vec{a}\cdot\vec{b}$ pour $\vec{a}\coord{4}{-2}$ et $\vec{b}\coord{1}{2}$.
        Que peut-on dire de ces vecteurs ?
  \item On donne $E(2\,;-1)$, $F(5\,;1)$ et $G(0\,;3)$. Calculer $\vect{EF}\cdot\vect{EG}$.
\end{questions}

\end{methodebox}

\begin{methodebox}[Méthode 12 --- Démontrer que deux droites sont perpendiculaires]
Les droites $(AB)$ et $(CD)$ sont perpendiculaires quand $\vect{AB}\cdot\vect{CD}=0$.
Pour deux droites données par une équation $ax+by+c=0$, on fait le produit scalaire de
leurs vecteurs directeurs $\vec{u}\coord{-b}{a}$ (méthode 6).
\etapeexemples
Le repère est orthonormé.
\begin{questions}[label=\alph*)]
  \item On donne $A(-2\,;1)$, $B(2\,;3)$, $C(3\,;1)$ et $D(1\,;5)$. Les droites $(AB)$
        et $(CD)$ sont-elles perpendiculaires ?
  \item On donne $E(1\,;0)$, $F(3\,;4)$, $G(-1\,;2)$ et $H(5\,;0)$. Les droites $(EF)$
        et $(GH)$ sont-elles perpendiculaires ?
  \item On donne $K(1\,;1)$, $L(4\,;2)$ et $M(0\,;4)$. Le triangle $KLM$ est-il
        rectangle en $K$ ?
  \item Les droites $(d_{1})$ : $3x+y-2=0$ et $(d_{2})$ : $x-3y+4=0$ sont-elles
        perpendiculaires ?
\end{questions}

\etape{À vous de jouer}
Le repère est orthonormé.
\begin{questions}[label=\alph*)]
  \item On donne $A(-3\,;0)$, $B(1\,;2)$, $C(2\,;-1)$ et $D(0\,;3)$. Les droites $(AB)$
        et $(CD)$ sont-elles perpendiculaires ?
  \item On donne $E(0\,;1)$, $F(2\,;5)$, $G(-2\,;3)$ et $H(4\,;1)$. Les droites $(EF)$
        et $(GH)$ sont-elles perpendiculaires ?
  \item On donne $R(0\,;2)$, $S(4\,;4)$ et $T(-1\,;4)$. Le triangle $RST$ est-il
        rectangle en $R$ ?
  \item Les droites $(d_{3})$ : $2x-5y+1=0$ et $(d_{4})$ : $5x+2y-3=0$ sont-elles
        perpendiculaires ?
\end{questions}

\end{methodebox}

\afaire{exercices 33 à 37}

\partiecours{Propriétés du produit scalaire}

\begin{proprietebox}[Propriété 10 --- Règles de calcul]
Pour tous vecteurs $\vec{u}$, $\vec{v}$, $\vec{w}$ et tout nombre réel $k$, le produit
scalaire se calcule comme un produit de nombres :
\begin{itemize}
  \item $\vec{u}\cdot\vec{v}=\vec{v}\cdot\vec{u}$ (l'ordre ne compte pas) ;
  \item $\vec{u}\cdot(\vec{v}+\vec{w})=\vec{u}\cdot\vec{v}+\vec{u}\cdot\vec{w}$
        (on développe) ;
  \item $(k\vec{u})\cdot\vec{v}=k\times(\vec{u}\cdot\vec{v})$ (les nombres sortent).
\end{itemize}
On en tire les identités remarquables :
\begin{gather*}
  \norme{\vec{u}+\vec{v}}^{2}=\norme{\vec{u}}^{2}+2\,\vec{u}\cdot\vec{v}+\norme{\vec{v}}^{2}
  \qquad
  \norme{\vec{u}-\vec{v}}^{2}=\norme{\vec{u}}^{2}-2\,\vec{u}\cdot\vec{v}+\norme{\vec{v}}^{2}\\
  (\vec{u}+\vec{v})\cdot(\vec{u}-\vec{v})=\norme{\vec{u}}^{2}-\norme{\vec{v}}^{2}.
\end{gather*}
Un vecteur retourné change de signe : $\vect{CA}=-\vect{AC}$, donc
$\vect{AB}\cdot\vect{CA}=-\vect{AB}\cdot\vect{AC}$.
\end{proprietebox}

\begin{methodebox}[Méthode 13 --- Calculer avec les règles du produit scalaire]
On développe comme avec des nombres, puis on remplace $\vec{u}\cdot\vec{u}$ par
$\norme{\vec{u}}^{2}$.
\etapeexemples
On donne $\norme{\vec{u}}=3$, $\norme{\vec{v}}=2$ et $\vec{u}\cdot\vec{v}=-1$. Calculer :
\begin{questions}[label=\alph*)]
  \item $\vec{u}\cdot(\vec{u}+\vec{v})$ ;
  \item $(\vec{u}+\vec{v})\cdot(\vec{u}-\vec{v})$ ;
  \item $(2\vec{u})\cdot(3\vec{v})$ ;
  \item $\norme{\vec{u}+\vec{v}}^{2}$, puis $\norme{\vec{u}+\vec{v}}$.
  \item On sait que $\vect{AB}\cdot\vect{AC}=5$. Calculer $\vect{AB}\cdot\vect{CA}$, puis
        $\vect{BA}\cdot\vect{CA}$.
\end{questions}

\etape{À vous de jouer}
On donne $\norme{\vec{u}}=4$, $\norme{\vec{v}}=1$ et $\vec{u}\cdot\vec{v}=2$. Calculer :
\begin{questions}[label=\alph*)]
  \item $\vec{u}\cdot(\vec{u}+\vec{v})$ ;
  \item $(\vec{u}+\vec{v})\cdot(\vec{u}-\vec{v})$ ;
  \item $(2\vec{u})\cdot(3\vec{v})$ ;
  \item $\norme{\vec{u}+\vec{v}}^{2}$, puis $\norme{\vec{u}+\vec{v}}$.
  \item On sait que $\vect{EF}\cdot\vect{EG}=7$. Calculer $\vect{EF}\cdot\vect{GE}$, puis
        $\vect{FE}\cdot\vect{GE}$.
\end{questions}

\end{methodebox}

\afaire{exercices 38 à 40}

\partiecours{Travail d'une force}

\begin{proprietebox}[Propriété 11 --- Travail d'une force]
\begin{minipage}[t]{0.58\linewidth}\raggedright
Un solide se déplace en ligne droite : son centre de gravité va de $A$ à $B$, sur l'axe
du déplacement. Le \textbf{travail} d'une force constante $\vec{F}$ pendant ce
déplacement est le produit scalaire
\[
  W=\vec{F}\cdot\vect{AB}=F\times AB\times\cos(\alpha),
\]
où $F$ est l'intensité de la force (en newtons, N), $AB$ la longueur du déplacement
(en mètres) et $\alpha$ l'angle entre la force et le déplacement. Le travail
s'exprime en joules (J).
\end{minipage}\hfill
\begin{minipage}[t]{0.38\linewidth}
\vspace{0pt}\centering
\begin{tikzpicture}[scale=0.75]
  \draw[fig,fill=black!6] (-0.7,-0.45) rectangle (0.7,0.45);
  \draw[aide] (4.3,-0.45) rectangle (5.7,0.45);
  \draw[axe] (-1.6,0) -- (6.8,0);
  \draw[deplacement] (0,0) -- (5,0);
  \node[below,font=\footnotesize] at (2.5,-0.5) {déplacement\,: $\vect{AB}$};
  \draw (1.5,0) arc (0:35:1.5);
  \node[font=\footnotesize,anchor=west] at (1.55,0.25) {$\alpha$};
  \draw[vecteur] (0,0) -- node[midway,above left,font=\footnotesize]{$\vec{F}$} ({3.2*cos(35)},{3.2*sin(35)});
  \pt{0,0}{below left}{A} \pt{5,0}{below right}{B}
\end{tikzpicture}
\end{minipage}
\end{proprietebox}

\begin{exemplebox}[Exemples --- trois forces sur une caisse tirée sur un sol horizontal]
\begin{minipage}[t]{0.58\linewidth}\raggedright
La caisse va de $A$ à $B$. Elle subit le frottement $\vec{f}$ du sol, opposé au
déplacement, son poids $\vec{P}$, vertical, et la traction $\vec{F}$ d'une corde, qui
fait un angle $\alpha$ avec le déplacement. Les forces sont dessinées à partir du centre
de gravité $G$.
\end{minipage}\hfill
\begin{minipage}[t]{0.38\linewidth}
\vspace{0pt}\centering
\troisforces{35}{$\alpha$}{$\vec{F}$}{$\vec{P}$}{$\vec{f}$}{$\vect{AB}$}
\end{minipage}

\cas{Force opposée au déplacement : le frottement $\vec{f}$}
\begin{calculs}
W_{f} &= f\times AB\times\cos(180^{\circ}) &\qquad& \rem{$\vec{f}$ et $\vect{AB}$ de sens contraires : angle de $180^{\circ}$}\\
W_{f} &= f\times AB\times(-1) && \rem{$\cos(180^{\circ})=-1$}\\
W_{f} &= -f\times AB &&
\end{calculs}
\explique{Le travail est \textbf{négatif} : le frottement freine la caisse (travail
résistant).}{une force opposée au déplacement}
\cas{Force perpendiculaire au déplacement : le poids $\vec{P}$}
\begin{calculs}
W_{P} &= P\times AB\times\cos(90^{\circ}) &\qquad& \rem{poids vertical, déplacement horizontal}\\
W_{P} &= P\times AB\times 0 && \rem{$\cos(90^{\circ})=0$}\\
W_{P} &= 0 &&
\end{calculs}
\explique{Le travail est \textbf{nul} : le poids ne fait ni avancer ni freiner la
caisse.}{une force perpendiculaire au déplacement}
\cas{Force quelconque : la traction $\vec{F}$}
\begin{calculs}
W_{F} &= F\times AB\times\cos(\alpha) &\qquad& \rem{angle $\alpha$ entre la force et le déplacement}\\
W_{F} &= AB\times\big(F\times\cos(\alpha)\big) && \rem{$F\times\cos(\alpha)$ : projection de $\vec{F}$ sur l'axe}
\end{calculs}
\explique{Seule la projection de $\vec{F}$ sur l'axe travaille. L'angle $\alpha$ est aigu,
donc $\cos(\alpha)>0$ : le travail est \textbf{positif}, la traction fait avancer la
caisse (travail moteur).}{une force inclinée sur le déplacement}
\end{exemplebox}

\begin{remarquebox}[Travail de la résultante]
Le travail de la résultante des forces est la \textbf{somme} des travaux de chaque
force :
$(\vec{F}+\vec{P}+\vec{f})\cdot\vect{AB}=\vec{F}\cdot\vect{AB}+\vec{P}\cdot\vect{AB}+\vec{f}\cdot\vect{AB}$
(on développe : propriété 10).
\end{remarquebox}

\begin{methodebox}[Méthode 14 --- Calculer le travail d'une ou de plusieurs forces]
\etapeexemples
\begin{minipage}[t]{0.58\linewidth}\raggedright
\begin{questions}[label=\alph*)]
  \item Un chariot avance de $12$~m sur un axe horizontal, tiré par une force
        $\vec{F}$ de $80$~N qui fait un angle de $30^{\circ}$ avec cet axe (schéma
        ci-contre). Calculer la projection de $\vec{F}$ sur l'axe, puis le travail de
        $\vec{F}$, au dixième.
\end{questions}
\end{minipage}\hfill
\begin{minipage}[t]{0.38\linewidth}
\vspace{0pt}\centering
\solidetire{30}{$\vec{F}$}{$80$~N}{$12$~m}
\end{minipage}

\smallskip
\begin{minipage}[t]{0.58\linewidth}\raggedright
Une luge avance de $20$~m sur un sol horizontal. Elle est tirée par une corde
($\vec{F_{1}}$ : $100$~N, à $30^{\circ}$ du sol) ; elle subit son poids ($\vec{P}$ :
$150$~N, vertical) et le frottement du sol ($\vec{f}$ : $40$~N, opposé au déplacement).
Calculer, au dixième :
\begin{questions}[label=\alph*),start=2]
  \item le travail du poids ;
  \item le travail de la traction $\vec{F_{1}}$ ;
  \item le travail du frottement ;
  \item le travail de la résultante des trois forces.
\end{questions}
\end{minipage}\hfill
\begin{minipage}[t]{0.38\linewidth}
\vspace{0pt}\centering
\troisforces{30}{$30^{\circ}$}{$\vec{F_{1}}$}{$\vec{P}$}{$\vec{f}$}{$20$~m}
\end{minipage}

\etape{À vous de jouer}
\begin{minipage}[t]{0.58\linewidth}\raggedright
\begin{questions}[label=\alph*)]
  \item Une luge avance de $5$~m sur un axe horizontal, tirée par une force
        $\vec{F}$ de $120$~N qui fait un angle de $40^{\circ}$ avec cet axe (schéma
        ci-contre). Calculer la projection de $\vec{F}$ sur l'axe, puis le travail de
        $\vec{F}$, au dixième.
\end{questions}
\end{minipage}\hfill
\begin{minipage}[t]{0.38\linewidth}
\vspace{0pt}\centering
\solidetire{40}{$\vec{F}$}{$120$~N}{$5$~m}
\end{minipage}

\smallskip
\begin{minipage}[t]{0.58\linewidth}\raggedright
Un chariot avance de $15$~m sur un sol horizontal. Il est tiré par un câble
($\vec{F_{1}}$ : $200$~N, à $20^{\circ}$ du sol) ; il subit son poids ($\vec{P}$ :
$300$~N, vertical) et le frottement du sol ($\vec{f}$ : $60$~N, opposé au déplacement).
Calculer, au dixième :
\begin{questions}[label=\alph*),start=2]
  \item le travail du poids ;
  \item le travail de la traction $\vec{F_{1}}$ ;
  \item le travail du frottement ;
  \item le travail de la résultante des trois forces.
\end{questions}
\end{minipage}\hfill
\begin{minipage}[t]{0.38\linewidth}
\vspace{0pt}\centering
\troisforces{20}{$20^{\circ}$}{$\vec{F_{1}}$}{$\vec{P}$}{$\vec{f}$}{$15$~m}
\end{minipage}

\end{methodebox}

\afaire{exercices 41 à 44}

\partiecours{Calcul d'un angle}

\begin{proprietebox}[Propriété 12 --- Cosinus d'un angle]
En écrivant le produit scalaire de deux façons (avec les coordonnées, puis avec
l'angle $\alpha$ entre deux vecteurs non nuls $\vec{u}$ et $\vec{v}$) :
\[
  \cos(\alpha)=\frac{\vec{u}\cdot\vec{v}}{\norme{\vec{u}}\times\norme{\vec{v}}}.
\]
La touche $\cos^{-1}$ (ou Arccos) de la calculatrice, en mode degré, donne ensuite
l'angle $\alpha$, entre $0^{\circ}$ et $180^{\circ}$.
\end{proprietebox}

\begin{methodebox}[Méthode 15 --- Calculer un angle avec le produit scalaire]
\etapeexemples
Le repère est orthonormé. Donner les angles au dixième de degré.
\begin{questions}[label=\alph*)]
  \item Calculer l'angle entre $\vec{u}\coord{3}{1}$ et $\vec{v}\coord{1}{2}$.
  \item Calculer l'angle entre $\vec{a}\coord{2}{1}$ et $\vec{b}\coord{-3}{1}$.
  \item On donne $A(1\,;1)$, $B(5\,;1)$ et $C(2\,;4)$. Calculer l'angle $\widehat{BAC}$.
\end{questions}

\etape{À vous de jouer}
Le repère est orthonormé. Donner les angles au dixième de degré.
\begin{questions}[label=\alph*)]
  \item Calculer l'angle entre $\vec{u}\coord{4}{2}$ et $\vec{v}\coord{1}{3}$.
  \item Calculer l'angle entre $\vec{a}\coord{1}{2}$ et $\vec{b}\coord{-3}{1}$.
  \item On donne $E(0\,;1)$, $F(3\,;1)$ et $G(1\,;3)$. Calculer l'angle $\widehat{FEG}$.
\end{questions}

\end{methodebox}

\afaire{exercices 45 à 47}

\partiecours{Théorème d'Al-Kashi}

\begin{proprietebox}[Propriété 13 --- Théorème d'Al-Kashi]
\begin{minipage}[t]{0.58\linewidth}\raggedright
Dans un triangle $ABC$, si l'on connaît deux côtés et l'angle entre eux, on calcule le
troisième côté :
\[
  BC^{2}=AB^{2}+AC^{2}-2\times AB\times AC\times\cos\big(\widehat{BAC}\big).
\]
C'est le théorème de Pythagore « corrigé » : si l'angle en $A$ est droit,
$\cos(90^{\circ})=0$ et on retrouve $BC^{2}=AB^{2}+AC^{2}$.
\end{minipage}\hfill
\begin{minipage}[t]{0.38\linewidth}
\vspace{0pt}\centering
\begin{tikzpicture}[scale=0.6]
  \draw[fig] (0,0) -- (5,0) -- (1.5,2.6) -- cycle;
  \draw[coulExo,line width=0.8pt] (0.7,0) arc (0:60:0.7);
  \node[font=\footnotesize,coulExo,anchor=west] at (0.75,0.5) {$\widehat{BAC}$};
  \pt{0,0}{below left}{A} \pt{5,0}{below right}{B} \pt{1.5,2.6}{above}{C}
\end{tikzpicture}
\end{minipage}

\smallskip
\demo{avec les règles de calcul (propriété 10) :}
\begin{calculs}
BC^{2} &= \norme{\vect{AC}-\vect{AB}}^{2} &\qquad& \rem{$\vect{BC}=\vect{BA}+\vect{AC}$}\\
BC^{2} &= AC^{2}-2\,\vect{AB}\cdot\vect{AC}+AB^{2} && \rem{identité remarquable}\\
BC^{2} &= AB^{2}+AC^{2}-2\times AB\times AC\times\cos\big(\widehat{BAC}\big) && \rem{définition 4}
\end{calculs}
\end{proprietebox}

\begin{methodebox}[Méthode 16 --- Calculer une longueur avec le théorème d'Al-Kashi]
\etapeexemples
\begin{questions}[label=\alph*)]
  \item Dans le triangle $ABC$, $AB=5$, $AC=8$ et $\widehat{BAC}=60^{\circ}$. Calculer $BC$.
  \item Un bras de robot est formé de deux segments articulés, de $3$~m et de $4$~m, qui
        font entre eux un angle de $120^{\circ}$. Calculer la distance entre les deux
        extrémités du bras, au centième.
\end{questions}

\etape{À vous de jouer}
\begin{questions}[label=\alph*)]
  \item Dans le triangle $EFG$, $EF=3$, $EG=8$ et $\widehat{FEG}=60^{\circ}$. Calculer $FG$.
  \item Un bras de grue est formé de deux segments articulés, de $2$~m et de $5$~m, qui
        font entre eux un angle de $120^{\circ}$. Calculer la distance entre les deux
        extrémités du bras, au centième.
\end{questions}

\end{methodebox}

\afaire{exercices 48 à 50}

\end{document}
